\documentclass[10pt]{extarticle}
\usepackage[a4paper,margin=1.5cm]{geometry}
\usepackage[T1]{fontenc}
\usepackage{amsmath,amssymb}
\usepackage{graphicx}
\usepackage[export]{adjustbox}
\usepackage{multicol}
\usepackage{enumitem}
\usepackage{booktabs,array,tabularx}
\setlength{\columnsep}{0.7cm}
\setlength{\parindent}{0pt}
\setlength{\parskip}{3pt}
\pagestyle{empty}
\makeatletter
\newcommand{\qfigdef}[4]{\@namedef{qf@file@#1}{#2}\@namedef{qf@w@#1}{#3}\@namedef{qf@wide@#1}{#4}}
\newcommand{\qf@check}[1]{\@ifundefined{qf@file@#1}{\PackageError{score2}{Unknown figure id #1}{}}{}}
\newcommand{\qf@img}[4][]{\includegraphics[width=#2\dimexpr\@nameuse{qf@w@#4}pt\relax,max width=#3,max height=0.45\textheight,#1]{\@nameuse{qf@file@#4}}}
\newif\ifqf@full
\newif\ifqf@boxed
\newsavebox{\qf@box}
\newcommand{\qf@open}{\global\qf@fulltrue\par\global\qf@boxedtrue
  \setbox\qf@box\vbox\bgroup\hsize\textwidth\linewidth\textwidth}
\newcommand{\qf@emit}[1]{\par\egroup\global\qf@boxedfalse
  \ifdim\dimexpr\ht\qf@box+\dp\qf@box\relax<#1\textheight\noindent\usebox{\qf@box}\par\else\unvbox\qf@box\fi}
\newcommand{\qfullsplit}{\ifqf@boxed\qf@emit{0.85}\fi}
\newcommand{\qf@close}{\ifqf@boxed\qf@emit{0.45}\fi\global\qf@fullfalse}
\newcommand{\qfullbegin}{\end{multicols}\qf@open}
\newcommand{\qfullend}{\qf@close\begin{multicols}{2}}
\newcommand{\qfullstart}{\qf@open}
\newcommand{\qfulldone}{\qf@close}
\newcommand{\qfig}[2][1]{\qf@check{#2}%
  \ifqf@full\qf@blk{#1}{#2}\else
  \ifnum\@nameuse{qf@wide@#2}=1 \expandafter\@firstoftwo\else\expandafter\@secondoftwo\fi
  {\end{multicols}\begin{center}\qf@img{#1}{\textwidth}{#2}\end{center}\begin{multicols}{2}}%
  {\qf@blk{#1}{#2}}\fi}
\newcommand{\qf@blk}[2]{\par\smallskip{\centering\qf@img{#1}{\linewidth}{#2}\par}\smallskip}
\newcommand{\qfigi}[2][1]{\qf@check{#2}\qf@img[valign=c]{#1}{\linewidth}{#2}}
\makeatother
\newcommand{\papersec}[1]{\par\bigskip{\centering\Large\bfseries PAPER-#1\par}\medskip}
\newcommand{\partsec}[2]{\par\medskip\noindent{\large\bfseries PART-#1: #2}\par\smallskip}
\newcommand{\sect}[1]{\par\smallskip\noindent\textbf{#1}\par}
\newcommand{\qn}[1]{\par\smallskip\noindent\textbf{#1.}~}
\newcommand{\paraq}[2]{\par\smallskip\noindent\textbf{Paragraph for Questions #1 and #2}\par}
\newenvironment{qoptions}{\begin{itemize}[leftmargin=1.8em,itemsep=1pt,topsep=2pt,label={}]}{\end{itemize}}
\newcommand{\opt}[2]{\item[(#1)] #2}
% figure registry (generated by score2_tools.py finalize)
\qfigdef{p1_phy_s1_q01_1}{Score-I - Test 08 (09 Feb 2022)-img/p1_phy_s1_q01_1.png}{330.6}{1}
\qfigdef{p1_phy_s2_q05_1}{Score-I - Test 08 (09 Feb 2022)-img/p1_phy_s2_q05_1.png}{126.4}{0}
\qfigdef{p1_phy_s2_q06_1}{Score-I - Test 08 (09 Feb 2022)-img/p1_phy_s2_q06_1.png}{113.1}{0}
\qfigdef{p1_phy_s2_q07_1}{Score-I - Test 08 (09 Feb 2022)-img/p1_phy_s2_q07_1.png}{260.4}{1}
\qfigdef{p1_phy_s3_q03_1}{Score-I - Test 08 (09 Feb 2022)-img/p1_phy_s3_q03_1.png}{99.2}{0}
\qfigdef{p1_phy_s4_q09_1}{Score-I - Test 08 (09 Feb 2022)-img/p1_phy_s4_q09_1.png}{184.6}{0}
\qfigdef{p1_chem_s1_q01_A}{Score-I - Test 08 (09 Feb 2022)-img/p1_chem_s1_q01_A.png}{108.3}{0}
\qfigdef{p1_chem_s1_q01_B}{Score-I - Test 08 (09 Feb 2022)-img/p1_chem_s1_q01_B.png}{108.3}{0}
\qfigdef{p1_chem_s1_q01_C}{Score-I - Test 08 (09 Feb 2022)-img/p1_chem_s1_q01_C.png}{108.3}{0}
\qfigdef{p1_chem_s1_q01_D}{Score-I - Test 08 (09 Feb 2022)-img/p1_chem_s1_q01_D.png}{108.5}{0}
\qfigdef{p1_chem_s1_q03_A}{Score-I - Test 08 (09 Feb 2022)-img/p1_chem_s1_q03_A.png}{28.2}{0}
\qfigdef{p1_chem_s1_q03_B}{Score-I - Test 08 (09 Feb 2022)-img/p1_chem_s1_q03_B.png}{29.1}{0}
\qfigdef{p1_chem_s1_q03_C}{Score-I - Test 08 (09 Feb 2022)-img/p1_chem_s1_q03_C.png}{41.3}{0}
\qfigdef{p1_chem_s1_q03_D}{Score-I - Test 08 (09 Feb 2022)-img/p1_chem_s1_q03_D.png}{39.9}{0}
\qfigdef{p1_chem_s1_q04_1}{Score-I - Test 08 (09 Feb 2022)-img/p1_chem_s1_q04_1.png}{294.4}{1}
\qfigdef{p1_chem_s1_q04_A}{Score-I - Test 08 (09 Feb 2022)-img/p1_chem_s1_q04_A.png}{55.8}{0}
\qfigdef{p1_chem_s1_q04_B}{Score-I - Test 08 (09 Feb 2022)-img/p1_chem_s1_q04_B.png}{49.9}{0}
\qfigdef{p1_chem_s1_q04_C}{Score-I - Test 08 (09 Feb 2022)-img/p1_chem_s1_q04_C.png}{69.3}{0}
\qfigdef{p1_chem_s1_q04_D}{Score-I - Test 08 (09 Feb 2022)-img/p1_chem_s1_q04_D.png}{73.0}{0}
\qfigdef{p1_chem_s2_q09_A}{Score-I - Test 08 (09 Feb 2022)-img/p1_chem_s2_q09_A.png}{111.5}{0}
\qfigdef{p1_chem_s2_q09_B}{Score-I - Test 08 (09 Feb 2022)-img/p1_chem_s2_q09_B.png}{131.3}{0}
\qfigdef{p1_chem_s2_q09_C}{Score-I - Test 08 (09 Feb 2022)-img/p1_chem_s2_q09_C.png}{101.1}{0}
\qfigdef{p1_chem_s2_q09_D}{Score-I - Test 08 (09 Feb 2022)-img/p1_chem_s2_q09_D.png}{107.8}{0}
\qfigdef{p1_chem_s2_q10_A}{Score-I - Test 08 (09 Feb 2022)-img/p1_chem_s2_q10_A.png}{113.2}{0}
\qfigdef{p1_chem_s2_q10_C}{Score-I - Test 08 (09 Feb 2022)-img/p1_chem_s2_q10_C.png}{92.6}{0}
\qfigdef{p1_chem_s2_q10_D}{Score-I - Test 08 (09 Feb 2022)-img/p1_chem_s2_q10_D.png}{108.2}{0}
\qfigdef{p1_chem_s3_q05_ref}{Score-I - Test 08 (09 Feb 2022)-img/p1_chem_s3_q05_ref.png}{41.4}{0}
\qfigdef{p1_chem_s3_q05_list}{Score-I - Test 08 (09 Feb 2022)-img/p1_chem_s3_q05_list.png}{224.3}{0}
\qfigdef{p1_chem_s3_q06_list}{Score-I - Test 08 (09 Feb 2022)-img/p1_chem_s3_q06_list.png}{224.9}{0}
\qfigdef{p1_chem_s4_q08_2}{Score-I - Test 08 (09 Feb 2022)-img/p1_chem_s4_q08_2.png}{160.2}{0}
\qfigdef{p2_phy_s1_q03_1}{Score-I - Test 08 (09 Feb 2022)-img/p2_phy_s1_q03_1.png}{54.2}{0}
\qfigdef{p2_phy_s1_q04_1}{Score-I - Test 08 (09 Feb 2022)-img/p2_phy_s1_q04_1.png}{245.2}{0}
\qfigdef{p2_phy_s2_para7_8_1}{Score-I - Test 08 (09 Feb 2022)-img/p2_phy_s2_para7_8_1.png}{124.9}{0}
\qfigdef{p2_phy_s2_para9_10_1}{Score-I - Test 08 (09 Feb 2022)-img/p2_phy_s2_para9_10_1.png}{142.7}{0}
\qfigdef{p2_phy_s3_para1_2_1}{Score-I - Test 08 (09 Feb 2022)-img/p2_phy_s3_para1_2_1.png}{147.9}{0}
\qfigdef{p2_phy_s3_para3_4_1}{Score-I - Test 08 (09 Feb 2022)-img/p2_phy_s3_para3_4_1.png}{174.6}{0}
\qfigdef{p2_phy_s3_para5_6_1}{Score-I - Test 08 (09 Feb 2022)-img/p2_phy_s3_para5_6_1.png}{102.4}{0}
\qfigdef{p2_phy_s4_q07_1}{Score-I - Test 08 (09 Feb 2022)-img/p2_phy_s4_q07_1.png}{171.9}{0}
\qfigdef{p2_phy_s4_q08_1}{Score-I - Test 08 (09 Feb 2022)-img/p2_phy_s4_q08_1.png}{76.5}{0}
\qfigdef{p2_phy_s4_q09_1}{Score-I - Test 08 (09 Feb 2022)-img/p2_phy_s4_q09_1.png}{199.3}{0}
\qfigdef{p2_chem_s1_q05_A}{Score-I - Test 08 (09 Feb 2022)-img/p2_chem_s1_q05_A.png}{150.1}{0}
\qfigdef{p2_chem_s1_q05_B}{Score-I - Test 08 (09 Feb 2022)-img/p2_chem_s1_q05_B.png}{115.3}{0}
\qfigdef{p2_chem_s1_q05_C}{Score-I - Test 08 (09 Feb 2022)-img/p2_chem_s1_q05_C.png}{162.1}{0}
\qfigdef{p2_chem_s1_q05_D}{Score-I - Test 08 (09 Feb 2022)-img/p2_chem_s1_q05_D.png}{160.9}{0}
\qfigdef{p2_chem_s1_q06_A}{Score-I - Test 08 (09 Feb 2022)-img/p2_chem_s1_q06_A.png}{130.8}{0}
\qfigdef{p2_chem_s1_q06_B}{Score-I - Test 08 (09 Feb 2022)-img/p2_chem_s1_q06_B.png}{134.2}{0}
\qfigdef{p2_chem_s1_q06_C}{Score-I - Test 08 (09 Feb 2022)-img/p2_chem_s1_q06_C.png}{134.4}{0}
\qfigdef{p2_chem_s1_q06_D}{Score-I - Test 08 (09 Feb 2022)-img/p2_chem_s1_q06_D.png}{118.2}{0}
\qfigdef{p2_chem_s4_q09_1}{Score-I - Test 08 (09 Feb 2022)-img/p2_chem_s4_q09_1.png}{324.9}{1}

\begin{document}
\begin{center}{\Large\bfseries JEE (Main+Advanced) Score-I --- Test 08 (09 Feb 2022)}\end{center}
\begin{multicols}{2}
\raggedcolumns
\qfullbegin
\papersec{1}
\partsec{1}{PHYSICS}
\sect{Section-I(i)}
\qn{1}
\qfig{p1_phy_s1_q01_1}
\qfullsplit
In the picture above, 50 vernier scale divisions coincide with 49 main scale divisions which are 1 mm each. What is the reading on scale?
\begin{qoptions}
\opt{A}{37.40 mm}
\opt{B}{37.41 mm}
\opt{C}{37.49 mm}
\opt{D}{37.46 mm}
\end{qoptions}
\qfullend
\qn{2} A Lightweight thin wire ring hanging on a soap film that is held in a circular frame. The mass of the ring is $m$, its radius $R$, the coefficient of surface tension of the film is $\sigma$, the diameter of the frame $D>2R$. The frame and ring are horizontal, their centres are on the same vertical line. What is the distance from the plane of the ring to the plane of the frame? Mass of the film can be neglected in comparison with the mass of the ring, Also $mg\ll\sigma R$.
\begin{qoptions}
\opt{A}{$h=\dfrac{mg}{4\pi\sigma}\ln\dfrac{D}{2R}$}
\opt{B}{$h=\dfrac{mg}{4\pi\sigma R}\left(\dfrac{D}{2}-R\right)$}
\opt{C}{$h=\dfrac{mg}{2\pi\sigma R}\left(\dfrac{D}{2}-R\right)$}
\opt{D}{$h=\dfrac{mg}{4\pi\sigma}\tan^{-1}\left(\dfrac{D}{2R}\right)$}
\end{qoptions}
\qn{3} A parallel plate capacitor with a capacity $C$ with an air dielectric consists of two large plates, located very close to each other. One of Plates is uncharged, another has a charge $Q$. We connect the plates with a conductor having a high resistance $R$. Estimate the amount of heat that will be dissipated in the conductor.
\begin{qoptions}
\opt{A}{$\dfrac{Q^2}{2C}$}
\opt{B}{$\dfrac{Q^2}{4C}$}
\opt{C}{$\dfrac{Q^2}{6C}$}
\opt{D}{$\dfrac{Q^2}{8C}$}
\end{qoptions}
\qn{4} In a YDSE apparatus, the intensity at the central maxima is 4 W/m$^2$. Both the slits are of equal width and the wavelength of the light used is 600 nm. An extremely thin glass plate of thickness 200 nm and refractive index 1.5 is placed in front of one of the slits and it is seen that the intensity at the centre of the screen (where earlier there was central maxima) is now 1.75 W/m$^2$. It is suspected that the glass plate absorbs some light. What is the new intensity at the central maxima ?
\begin{qoptions}
\opt{A}{3.00 W/m$^2$}
\opt{B}{2.25 W/m$^2$}
\opt{C}{3.60 W/m$^2$}
\opt{D}{2.40 W/m$^2$}
\end{qoptions}
\sect{Section-I(ii)}
\qn{5} A wire 30.0 cm long is held parallel to and 80.0 cm above a long wire carrying 200 A and resting on the floor (Fig.). The 30.0-cm wire is released at $t=0$ and falls, remaining parallel with the current-carrying wire as it falls. Assume that the falling wire accelerates at 10 m/s$^2$ (The wire is falling freely).
\qfig{p1_phy_s2_q05_1}
\begin{qoptions}
\opt{A}{The left end of the wire will be at a higher potential as it falls.}
\opt{B}{The emf is maximum when it is about to hit the wire}
\opt{C}{The emf at $t=0.2$ sec is 0.04 mV}
\opt{D}{The emf is maximum at the initial moment.}
\end{qoptions}
\qn{6} Consider an ohmmeter used to measure the resistance $R$ connected to it. It is basically an ammeter whose scale is calibrated to measure the resistance $R$. The internal resistance of the battery is $r$, its emf is $E$, and the internal resistance of the ammeter is $r_1$. The least count of the ammeter is $\Delta I$.
\qfig{p1_phy_s2_q06_1}
\begin{qoptions}
\opt{A}{The error in the measurement of $R$ is minimum when current is maximum}
\opt{B}{The error in the measurement of $R$ is maximum when current is maximum}
\opt{C}{The ohmmeter scale is linear (the divisions are equally spaced)}
\opt{D}{The sensitivity of the ohmmeter $\left(\dfrac{dI}{dR}\right)$ is maximum when current is maximum}
\end{qoptions}
\qfullbegin
\qn{7} They say that in Snell's archive they found an optical scheme, which depicted lens, Real object and its real image. The object and the image are of the same size and shape (but longitudinally inverted), and the longer line of the object lies on the optical axis of the lens. The image and object both are shown here. The image can be on either left hand side or right hand side of object. Focal length of the lens $=f$
\qfig{p1_phy_s2_q07_1}
\qfullsplit
\begin{qoptions}
\opt{A}{$O$ should be at a distance of $2f$ from the lens.}
\opt{B}{The lens should be on the right-hand side of the object}
\opt{C}{The focal length of the lens is $2a$ where $a$ is length of one square grid}
\opt{D}{The lens is converging in nature}
\end{qoptions}
\qfullend
\qn{8} A progressive transverse sine wave is passing through a string held along $x$ axis. Consider a small element of the string. Choose the correct statement(s):
\begin{qoptions}
\opt{A}{The element executes SHM in a direction perpendicular to $x$ axis.}
\opt{B}{When the element reaches the extreme position, its elastic potential energy is maximum and its kinetic energy is zero.}
\opt{C}{When the element reaches the mean position, its kinetic energy is maximum but its elastic potential energy is zero.}
\opt{D}{At any given instant of time, the kinetic energy and elastic potential energy of the element are the same.}
\end{qoptions}
\qn{9} A mole of helium in the process of expansion has constant molar heat capacity during the process which is equal to $C=15$ J / (mol - K). The gas does 20 J work during the process. Take $R=25/3$ J/mol K
\begin{qoptions}
\opt{A}{The heat given to the gas is 120 J}
\opt{B}{The temperature of the gas increases by 8 K}
\opt{C}{Heat extracted from the gas is 80 J}
\opt{D}{The pressure of the gas increases during the process.}
\end{qoptions}
\qn{10} An object of mass 2 kg is left to move along a horizontal surface, on which friction is not negligible. At one moment its kinetic energy is 900 J, and 2 s later it is only 400 J, while the object covers a distance of $\ell$.
\begin{qoptions}
\opt{A}{The object stops after covering a further distance of 40 m after covering the initial distance of $\ell$.}
\opt{B}{The object takes 4 sec more to stop after covering the distance of $\ell$.}
\opt{C}{The friction coefficient between the object and the surface is 1}
\opt{D}{$\ell=50$ m}
\end{qoptions}
\sect{Section-II(i)}
\paraq{1}{2}
Recently there has been a lot of hue and cry over the possibility of an asteroid hitting the earth. NASA reports that it has spotted an asteroid Bennu which may possibly hit the earth somewhere in 22$^{\text{nd}}$ century. The probability of its hitting the earth's surface is only one in 2000, yet social media sites are expert in spreading rumour about its possible impact. Assume that it was at rest at a large distance from earth and it is moving only under the influence of earth's gravitational field. Neglect the influence of any other celestial body like sun, moon etc. on the asteroid or the earth. The size of Bennu is much smaller than earth. (Take : escape speed from earth's surface as 11 km/s.)
\qn{1} If the asteroid were to hit the surface of the earth normally at the equator and gets embedded in the earth just below the surface, the duration of the day would become 30 hrs. What is the mass of the asteroid/mass of earth ?
\qn{2} What is the velocity $v$ (in km/s) of asteroid just before hitting the earth? Fill the value of $v^2$.
\paraq{3}{4}
Consider a small disc of mass $m$ and charge $+Q$ on the surface of a rough inclined plane. A uniform magnetic field exists perpendicular to the inclined plane as shown. The inclined plane is inclined at an angle of $\alpha$ to the horizontal. The inclined plane is very long and wide. The coefficient of friction is $\mu$. Take $\alpha=60^\circ$, $m=80$ gm, $q=100$ mC, $B=0.4$ T.
\qn{3} If $\mu=\dfrac{(\tan\alpha)}{2}$, what is the speed of the disc (in m/s) after a long time? The disc is released from rest.
\qfig{p1_phy_s3_q03_1}
\qn{4} Now assume that the width of the inclined plane is $2L$. At the initial moment, the disc is at a distance of $L$ from the left edge. Also assume that the friction coefficient is $\mu=\tan\alpha$. What can be the minimum initial speed along the inclined plane (in cm/s) so that the disc leaves the inclined plane from the edge? Take : $L=0.3$ m.
\paraq{5}{6}
One important application of radioactivity is the determination of the age of archaeological or geological samples. If an object contains radioactive nuclei when it is formed, then the decay of these nuclei marks the passage of time like a clock, half of the nuclei disintegrating during each half-life. If the half-life is known, a measurement of the number of nuclei present today relative to the number present initially can give the age of the sample. The radiocarbon technique utilizes the $^{14}_{\ 6}$C isotope of carbon, which undergoes $\beta^-$ decay with a half-life of $2\times10^{11}$ sec. This isotope is present in the earth's atmosphere at an equilibrium concentration of about one atom for every $7\times10^{11}$ atoms of normal carbon $^{12}_{\ 6}$C. It is often assumed that this value has remained constant over the years because $^{14}_{\ 6}$C is created when cosmic rays interact with the earth's upper atmosphere, a production method that offsets the loss via $\beta^-$ decay. Moreover, nearly all living organisms ingest the equilibrium concentration of $^{14}_{\ 6}$C. However, once an organism dies, metabolism no longer sustains the input of $^{14}_{\ 6}$C, and $\beta$ decay causes half of the $^{14}_{\ 6}$C nuclei to disintegrate every $2\times10^{11}$ sec. Take $\ln 2=0.7$
\qn{5} Find the activity of sample of blood of a living organism sample in Bq/gm of carbon atoms present in the sample.
\qn{6} In the Italian Alps, German tourists found a Stone-Age traveller, later dubbed the Ice Man, whose body had become trapped in a glacier. Material found with the body had a $^{14}_{\ 6}$C activity of about 0.125 Bq per gram of carbon. Find the age (in $10^{11}$ sec) of the Ice Man's remains.
\sect{Section-II(ii)}
\qn{7} Consider a dipole made up of charges 2.5 $\mu$C and $-2.5$ $\mu$C kept at (0.125 mm, 0, 0) and ($-0.125$mm, 0, 0) respectively. Find the magnitude of the electric field (in mN/C) due to this dipole at $(2\sqrt2\,\mathrm m,\ \sqrt{17}\,\mathrm m,\ 0)$.
\qn{8} A rod of cross-section $S=100$ cm$^2$ is held between two fixed solid walls 1m apart. The temperature of one wall is then raised, and a constant heat flux $Q=0.5$ W is established in the rod. What is the pressure (in $10^8$ Pa) exerted on the rod if the initial stress in the rod was zero? The thermal conductivity of the rod is $k=0.2$ W/m K, the linear expansion coefficient is $\alpha=2\times10^{-5}$/K, Young's modulus is $E=2\times10^{11}$ N/m$^2$.
\qn{9} A horizontal plate is in equilibrium above a source of water jet at a height $h$ as shown in the figure. The water splashes parallel to the plate and does not rebound. The gravitational force on the plate is balanced by the force exerted by the jet of water. If the plate is displaced by a small distance in the vertically downward or upward direction, it is seen to execute SHM with a time period $T$. Find $T$ (in sec.). The speed of jet when it strikes the plate is $\dfrac{40}{\pi}$ m/s.
\qfig{p1_phy_s4_q09_1}
\raggedcolumns
\columnbreak
\partsec{2}{CHEMISTRY}
\sect{Section-I(i)}
\qn{1} The correct curve for 3p orbital is -
\begin{qoptions}
\opt{A}{\qfigi{p1_chem_s1_q01_A}}
\opt{B}{\qfigi{p1_chem_s1_q01_B}}
\opt{C}{\qfigi{p1_chem_s1_q01_C}}
\opt{D}{\qfigi{p1_chem_s1_q01_D}}
\end{qoptions}
\columnbreak
\qn{2} In which of the following pair EN of first element is greater than second element while EA of first element is less than second element.
\begin{qoptions}
\opt{A}{As, P}
\opt{B}{O, N}
\opt{C}{P, Si}
\opt{D}{S, Cl}
\end{qoptions}
\qn{3} Which of following compound doesn't react with NaHCO$_3$ giving CO$_2$ ?
\begin{qoptions}
\opt{A}{\qfigi{p1_chem_s1_q03_A}}
\opt{B}{\qfigi{p1_chem_s1_q03_B}}
\opt{C}{\qfigi{p1_chem_s1_q03_C}}
\opt{D}{\qfigi{p1_chem_s1_q03_D}}
\end{qoptions}
\qfullbegin
\qn{4} Major product ``Z'' obtained in given reaction is ?
\qfig{p1_chem_s1_q04_1}
\qfullsplit
\begin{qoptions}
\opt{A}{\qfigi{p1_chem_s1_q04_A}}
\opt{B}{\qfigi{p1_chem_s1_q04_B}}
\opt{C}{\qfigi{p1_chem_s1_q04_C}}
\opt{D}{\qfigi{p1_chem_s1_q04_D}}
\end{qoptions}
\qfullend
\sect{Section-I(ii)}
\qn{5} Irreversible process is :

$\mathrm{H_2O(\ell,\ 1\ atm,\ -10^\circ C)\ \rightarrow\ H_2O\ (S,\ 1\ atm,\ -10^\circ C)}$

Given at $0^\circ$C \quad $\Delta H_{\mathrm{fusion}}=6$ kJ/mole

at $-10^\circ$C \quad $\Delta H_{\mathrm{fusion}}=5.55$ kJ/mole

$C_p$(ice) $=30$ J/K-mole

$C_p$(water) $=75$ J/K-mole

$\ln\dfrac{273}{263}=0.04$

Pick out the incorrect statements :
\begin{qoptions}
\opt{A}{$\Delta S_{\mathrm{total}}$ for process $=-20.4$ J/K-mole}
\opt{B}{$\Delta S_{\mathrm{surr}}$ for process $=21.1$ J/K-mole}
\opt{C}{$\Delta S_{\mathrm{total}}$ for process $=0.92$ J/K-mole}
\opt{D}{$\Delta S_{\mathrm{total}}$ for process $=22$ J/K-mole}
\end{qoptions}
\qn{6} In which of the following colloidal systems the dispersion medium is liquid -
\begin{qoptions}
\opt{A}{sol}
\opt{B}{gel}
\opt{C}{emulsion}
\opt{D}{froath}
\end{qoptions}
\qn{7} Which of the following order is/are CORRECT for indicated property?
\begin{qoptions}
\opt{A}{Pd $<$ Pt (Maximum number of electrons having $\ell=2$)}
\opt{B}{1$^{\text{st}}$ period $<$ 2$^{\text{nd}}$ period $<$ 3$^{\text{rd}}$ period (Number of s-block elements)}
\opt{C}{C $<$ N $<$ O ($Z_{\mathrm{eff}}$ on last electron)}
\opt{D}{Cr $<$ Mn (unpaired electrons)}
\end{qoptions}
\qn{8} Which of the following is \textbf{correct} order(s) for indicated bond length -
\begin{qoptions}
\opt{A}{$d_{\mathrm{N-N}}$ in $\mathrm{N_2H_4}>d_{\mathrm{N-N}}$ in $\mathrm{N_2F_4}$}
\opt{B}{$d_{\mathrm{N-N}}$ in $\mathrm{NH_2{-}NH_2}>d_{\mathrm{N-N}}$ in $\overset{\oplus}{\mathrm N}\mathrm H_3-\overset{\oplus}{\mathrm N}\mathrm H_3$}
\opt{C}{$d_{\mathrm{Si-Cl}}$ in $\mathrm{SiClF_3}>d_{\mathrm{Si-Cl}}$ in $\mathrm{SiClH_3}$}
\opt{D}{$d_{\mathrm{C-C}}$ in $\mathrm{C_2H_6}>d_{\mathrm{C-C}}$ in $\mathrm{C_2F_6}$}
\end{qoptions}
\qn{9} Which of the following has $-$COOH group in major product :
\begin{qoptions}
\opt{A}{\qfigi{p1_chem_s2_q09_A}}
\opt{B}{\qfigi{p1_chem_s2_q09_B}}
\opt{C}{\qfigi{p1_chem_s2_q09_C}}
\opt{D}{\qfigi{p1_chem_s2_q09_D}}
\end{qoptions}
\qn{10} Select reaction(s) with correct major product :
\begin{qoptions}
\opt{A}{\qfigi{p1_chem_s2_q10_A}}
\opt{B}{$\mathrm{PhCH_2COOH}\xrightarrow[\mathrm{CaO}/\Delta]{\mathrm{NaOH}}\mathrm{PhCH_3}$}
\opt{C}{\qfigi{p1_chem_s2_q10_C}}
\opt{D}{\qfigi{p1_chem_s2_q10_D}}
\end{qoptions}
\sect{Section-II(i)}
\paraq{1}{2}
Enthalpy of combustion of a given compounds is defined as following. It is the change in enthalpy when 1 mole of this compound combines with required amount of O$_2$ to give product in their stable forms. The enthalpy of combustion is usually measured by placing a known mass of the compound in a closed steel cointainer (known as bomb-calorimeter)

Given the data below :

$\Delta H^\circ_{\mathrm{comb}}(\mathrm{CH_4},\mathrm g)=-900$ kJ/mole;\par
$\Delta H^\circ_{\mathrm{comb}}(\mathrm C,\text{graphite})=-400$ kJ/mole

$\Delta H^\circ_{\mathrm{comb}}(\mathrm{H_2},\mathrm g)=-300$ kJ/mole ;\par
$\Delta H^\circ_{\mathrm{atomization}}(\mathrm C,\text{graphite})=700$ kJ/mole

and $\Delta H^\circ_{\mathrm f}(\mathrm H,\mathrm g)=200$ kJ/mole
\qn{1} What should be the bond enthalpy of C -- H bond in kJ/mole.
\qn{2} What should be magnitude of enthalpy of formation of $\mathrm{CH_{4(g)}}$ in kJ / mol?
\paraq{3}{4}
Crystal field theory (CFT) is an electrostatic model which consider metal-ligand bond to be ionic arising purely from electrostatic interaction between metal cation and ligand. It provide information about splitting of of orbitals and on that basis colour, stability of complex can be predicted.
\qn{3} Consider following complex formation reaction.

(i) $\mathrm{Co^{3+}_{(aq.)}}+\mathrm{KCN_{(excess)}}\rightleftharpoons$ complex ion `P'

(ii) $\mathrm{Ni^{2+}_{(aq.)}}+\mathrm{KCN_{(excess)}}\rightleftharpoons$ complex ion `Q'

and find the value of $\dfrac{(x+y)}{3}$ for complex ion P and Q.

where

$x=$ Total number of electron(s) present in $t_{2g}$ orbitals of central metal cation of complex ion `P'.

$y=$ Total number of electron(s) present in $d_{z^2}$ orbital of central metal cation of complex ion `Q'.
\qn{4} How many of the following statements are correct for octahedral ammine complex of Fe(II) and Fe(III)?

(i) Both Fe(II) and Fe(III) form outer orbital complex with NH$_3$.

(ii) Both form inner orbital complex with NH$_3$

(iii) Fe(II) form outer while Fe(III) form inner orbitals complex with NH$_3$

(iv) Both are diamagnetic

(v) Both are paramagnetic
\paraq{5}{6}
There are certain compounds which do not give S$_\mathrm{N}$1 and S$_\mathrm{N}$2 reaction. and rate of S$_\mathrm{N}$1 or S$_\mathrm{N}$2 $\propto$ stability of T.S. of RDS step.
\qn{5} How many compounds are more reactive than \qfigi{p1_chem_s3_q05_ref} towards S$_\mathrm{N}$1 (hydrolysis)?
\qfig{p1_chem_s3_q05_list}
\qn{6} How many compounds do not give ether as major product by Williumson ether synthesis in presence of $\overset{\oplus}{\mathrm{Na}}\overset{\ominus}{\mathrm O}\mathrm{CH_3}$ (sodium methoxide)
\qfig{p1_chem_s3_q06_list}
\sect{Section-II(ii)}
\qn{7} 0.01 moles of sulphuryl chloride $\mathrm{SO_2Cl_2(g)}$ is taken in a sealed tube \& heated to 400 K. It decomposes following 1$^{\text{st}}$ order kinetics according to reaction:

$\mathrm{SO_2Cl_2(g)\rightarrow SO_2(g)+Cl_2(g)}$

The tube is broken after 4 hours \& gas passed through a 25 ml of an acidified 1 N iodine solution where all SO$_2$ is oxidised to $\mathrm{SO_4^{2-}}$. The resulting solution required 10 ml of 1 M Hypo solution. What is half life in hours.
\qn{8} $\mathrm{P_4+SOCl_2}\xrightarrow{\ \Delta\ }$ Products

Find out total number of nonplanar and polar molecules in balanced equation for one mole of $\mathrm{P_4}$.
\qn{9}
\qfig{p1_chem_s4_q08_2}
Number of $\alpha$-Hydrogen atoms in Hydrocarbon formed as major product.
\raggedcolumns
\partsec{3}{MATHEMATICS}
\sect{Section-I(i)}
\qn{1} The solution of differential equation $y_1y_3=3y_2^{\,2}$ is (where $a,b,c$ are arbitrary constants \& $y_n$ denotes $n^{\text{th}}$ derivative of $y$ w r t `$x$')
\begin{qoptions}
\opt{A}{$x=ay^2+by+c$}
\opt{B}{$y=ax^2+bx+c$}
\opt{C}{$y=ae^x+be^{-x}+c$}
\opt{D}{$x=ae^y+be^{-y}+c$}
\end{qoptions}
\qn{2} If value of product of roots of equation $x^2-2ax+8-a=0$ lies between its roots, then maximum possible integral value of product of roots is
\begin{qoptions}
\opt{A}{1}
\opt{B}{4}
\opt{C}{6}
\opt{D}{7}
\end{qoptions}
\qn{3} $PQ$ and $PR$ are tangents to parabola $y^2=8x$ where $P$ lies on $x$-axis and $Q(2,4)$ and $R$ are points on the parabola. If normal to parabola at $R$ intersects the parabola again at $S$, then point of intersection of tangents at $Q$ and $S$ is
\begin{qoptions}
\opt{A}{$(-6,4)$}
\opt{B}{$(6,8)$}
\opt{C}{$(3,4)$}
\opt{D}{$(8,6)$}
\end{qoptions}
\qn{4} If $\hat a$, $\hat b$ and $\hat c$ are unit vectors such that $\hat a\wedge\hat b=\hat b\wedge\hat c=\hat c\wedge\hat a=60^\circ$ then $(\hat a\times\hat b)\times(\hat a+\hat b+\hat c)=\lambda\left(\vec b-\vec a\right)$, where $\lambda$ is
\begin{qoptions}
\opt{A}{$-2$}
\opt{B}{2}
\opt{C}{1}
\opt{D}{$-1$}
\end{qoptions}
\sect{Section-I(ii)}
\qn{5} If $R^3-5rR^2+8r^2R-4r^3=0$ in a $\Delta ABC$, then $\Delta ABC$ will be (all symbols are in usual notation) -
\begin{qoptions}
\opt{A}{equilateral}
\opt{B}{scalene}
\opt{C}{right angled}
\opt{D}{acute angled}
\end{qoptions}
\qn{6} Given that $1+2|z|^2=|z^2+1|^2+2|z+1|^2$, then
\begin{qoptions}
\opt{A}{$|z|=1$}
\opt{B}{$|z|=2$}
\opt{C}{$|z(z+1)|=1$}
\opt{D}{$|z(z+1)|=2$}
\end{qoptions}
\qn{7} Let $f$ be a differentiable function $\forall\ x\in R$ where $f(1)=1$ \& $f(3)=4$ \& $f'(x)\ge1\ \forall\ x\in R$, then $f(2)$ CANNOT be equal to
\begin{qoptions}
\opt{A}{1.5}
\opt{B}{2}
\opt{C}{2.5}
\opt{D}{3}
\end{qoptions}
\qn{8} Let $f:(0,\infty)\to R$ and $F(x)=\displaystyle\int_0^x t\,f(t)\,dt,\ \forall\ x>0$. If $F(x^2)=x^4+x^5,\ \forall\ x>0$ then -
\begin{qoptions}
\opt{A}{$f(x)$ is onto function}
\opt{B}{$f(4)=7$}
\opt{C}{$f(x)$ is continuous every where}
\opt{D}{$\displaystyle\int_0^4 f(x)\,dx=\dfrac{64}{3}$}
\end{qoptions}
\qn{9} $OAB$ is an isosceles triangle with $OA=OB$. If $O$ is origin, $A$ is $(-20,-6,-5)$ and equation of base $AB$ is $\dfrac{x+16}{2}=\dfrac{y}{3}=\dfrac{z+1}{2}$, then -
\begin{qoptions}
\opt{A}{If length of altitude from $O$ to the base $AB$ is $3\sqrt N$ then $N$ is 21.}
\opt{B}{If length of altitude from $O$ to the base $AB$ is $3\sqrt N$ then $N$ is 18.}
\opt{C}{If $B$ has co-ordinates $(a,b,c)$ then $a+b+c$ is 25}
\opt{D}{If $B$ has co-ordinates $(a,b,c)$ then $a+b+c$ is 30}
\end{qoptions}
\qn{10} If $A$ is an invertible square matrix of order less than 4 such that $|A-A^T|\ne0$ and $B=\operatorname{adj}A$, then which of the following is/are always true -
\begin{qoptions}
\opt{A}{$\operatorname{adj}(B^2A^{-1}B^{-1}A)=A$}
\opt{B}{$\operatorname{adj}(B^2A^{-1}B^{-1}A)=B$}
\opt{C}{Order of $A$ is 2}
\opt{D}{Order of $A$ is 3}
\end{qoptions}
\sect{Section-II(i)}
\paraq{1}{2}
Given $f(x)=\begin{cases}(x+1)^2 & -3<x\le0\\[2pt] 2x^{1/3}+a & 0<x\le1\\[2pt] 2-x & 1<x<2\end{cases}$ has maxima at $x=0,1$.
\qn{1} Number of possible integral values of $a$ is
\qn{2} Total number of extrema of $f(x)$ is
\paraq{3}{4}
14 chairs are arranged in a line and 6 of them are to be occupied by 6 persons (one on each chair)
\qn{3} Number of such arrangements in which no two persons are side by side is $^9C_k\cdot k!$, then $k$ is
\qn{4} Number of such arrangements in which each person has exactly one neighbour is $^9C_k\cdot k!$, then $k$ is
\paraq{5}{6}
Consider a family of circles whose equation is $x^2+y^2-6x+y(\lambda-4)+9-2\lambda=0$. Let this family passes through two fixed point $A(x_1,y_1)$ and $B(x_2,y_2)$.
\qn{5} The value of $(x_1+x_2+y_1+y_2)$ is equal to
\qn{6} If $S_2=0$ and $S_3=0$ be members of given family having equal radii and cutting each other orthogonally. Then the radius of these circles is
\sect{Section-II(ii)}
\qn{7} 3 natural numbers $x,y,z$ are in G.P. \& $x,y^2,z$ are in A.P., then the minimum value of $x+2y+z$ is equal to
\qn{8} A chord $AB$ of an ellipse $\dfrac{x^2}{9}+\dfrac{y^2}{\lambda^2}=1$ subtends a right angle at the centre of ellipse $O$ (where $O$ is origin), also $\dfrac{1}{OA^2}+\dfrac{1}{OB^2}=\dfrac{14}{45}$, then $\lambda^2$ is equal to
\qn{9} Four fair dice are thrown simultaneously. If the probability that the highest number obtained is 4 is $\dfrac{25a}{1296}$ then `$a$' is equal to
\raggedcolumns
\papersec{2}
\partsec{1}{PHYSICS}
\sect{Section-I(i)}
\qn{1} Two concentric conducting spheres are uncharged, and in the space between them, at a distance $L$ from the centre, a point charge $Q$ is fixed. The radius of the inner sphere is $r$ and the radius of the outer sphere is $R$. Now both spheres are connected by a thin conductor.
\begin{qoptions}
\opt{A}{The potential difference between spheres in the initial state is $kQ\left(\dfrac1L-\dfrac1R\right)$}
\opt{B}{The potential difference between the spheres in the initial state is $kQ\left(\dfrac1r-\dfrac1L\right)$}
\opt{C}{The final common potential of both the spheres is $k\left(\dfrac QR\right)$}
\opt{D}{A charge of $Q\,\dfrac{\frac1L-\frac1R}{\frac1r-\frac1R}$ will flow through the conductor.}
\end{qoptions}
\qn{2} Ram performs optical bench experiment to find out the focal length of a concave lens. For this purpose, he first takes a convex lens and measures the position of object needle and image needle. Then he combines convex lens and concave lens and again measures the object needle and image needle position. The optical bench has no bench error. The readings obtained by Ram in the two experiments are shown in the table below.
\begin{center}\small\setlength{\tabcolsep}{3pt}
\begin{tabular}{|p{1.9cm}|c|c|c|}\hline
 & Object needle & Lens & Image needle\\ \hline
Convex lens & 37.1 cm & 50.6 cm & 77.6 cm\\ \hline
Convex and Concave lens & 37.1 cm & 50.6 cm & 91.1 cm\\ \hline
\end{tabular}\end{center}
\begin{qoptions}
\opt{A}{The focal length of the convex lens is 9.0 cm.}
\opt{B}{The focal length of the concave lens is 81 cm in modulus.}
\opt{C}{The error in focal length of convex lens is approximately 0.1 cm.}
\opt{D}{The error in focal length of concave lens is approximately 1.0 cm.}
\end{qoptions}
\qn{3} In the figure shown below, cross represents an incandescent light bulb which is connected across a rheostat as shown. The battery has some internal resistance. As the jockey in the rheostat is moved from left to right
\qfig{p2_phy_s1_q03_1}
\begin{qoptions}
\opt{A}{The brightness of the bulb increases.}
\opt{B}{The current in the rheostat decreases.}
\opt{C}{The current through the battery increases and then decreases.}
\opt{D}{The current through the battery decreases and then increases.}
\end{qoptions}
\qn{4} Two bodies of equal mass are heated at the same uniform rate under identical conditions. Their temperatures are shown in the following figure.
\qfig{p2_phy_s1_q04_1}
\begin{qoptions}
\opt{A}{Melting points of A and B are $40^\circ$ C and $60^\circ$ C respectively.}
\opt{B}{The ratio of their latent heat of fusion $L_a/L_b$ is 0.75.}
\opt{C}{Their ratio of specific heats in solid state $S_a/S_b$ is 1 : 3}
\opt{D}{Their ratio of specific heats in liquid state $S_a/S_b$ is 1 : 2}
\end{qoptions}
\qn{5} A tube is open at one end and closed at the other by a vibration diaphragm which may be assumed to be a pressure node. It is found that when the frequency of the diaphragm is 2000Hz, a stationary wave pattern is set up in the tube and the distance between adjacent nodes is then 8.0 cm. When the frequency is gradually reduced the stationary wave pattern disappears but another stationary wave pattern reappears at a frequency of 1600 Hz. Neglect end correction.
\begin{qoptions}
\opt{A}{The speed of sound in air is 320 m/s.}
\opt{B}{The distance between adjacent nodes at a frequency of 1600 Hz is 0.2 m}
\opt{C}{The length of the tube is 40 cm}
\opt{D}{The next lower frequency at which a stationary wave pattern will be obtained is 800 Hz.}
\end{qoptions}
\qn{6} A small hole of radius 1 cm is made in a blackened spherical shell of radius 2m. Inside the shell we have another concentric blackened sphere of radius 1m. The inner sphere is perfectly conducting. Sun light of intensity 1600 W/m$^2$ falls through the hole on the inner shell in a direction perpendicular to the surface. The shell is also perfectly conducting and is located in vacuum. No radiations are falling from sun on the outer surface of the shell. Choose the correct statement(s) : Take : $\sigma=\dfrac{17}{3}\times10^{-8}\ \mathrm{W/m^2K^4}$
\begin{qoptions}
\opt{A}{The radiation spectrum coming out from the hole will be same as that of a black body at the temperature of the inner sphere.}
\opt{B}{The steady state temperature of the sphere is close to 30 K}
\opt{C}{The steady state temperature of the shell is close to 20 K}
\opt{D}{The steady state temperature of the shell is close to 10 K}
\end{qoptions}
\sect{Section-I(ii)}
\paraq{7}{8}
A fly $F$ is sitting on a glass slab $S$, 45cm thick \& of refractive index 3/2. The slab covers the top of a container $C$ containing water (refractive index 4/3) upto a height of 20 cm. The bottom of container is closed by a concave mirror $M$ of radius of curvature 40cm.
\qfig{p2_phy_s2_para7_8_1}
\qn{7} Locate the final image formed by all refractions \& reflection assuming paraxial rays.
\begin{qoptions}
\opt{A}{42.5 cm from the mirror}
\opt{B}{30 cm from the mirror}
\opt{C}{22.5 cm from the mirror}
\opt{D}{57.5 cm from the mirror}
\end{qoptions}
\qn{8} The length of the fly is 1mm. What is the length of its image after all refractions and reflection? The length of the fly is perpendicular to the principal axis of the mirror.
\begin{qoptions}
\opt{A}{2 mm}
\opt{B}{0.5 mm}
\opt{C}{0.33 mm}
\opt{D}{4 mm}
\end{qoptions}
\paraq{9}{10}
Two thin sheets on the surfaces of a slope of angle of elevation of $\theta=30^\circ$ form a pair of rails. The distance between the sheets is $d=1.6$ cm. A solid ball of diameter of 2 cm is placed to the rails and its rolls down without slipping. Assume that the ball does not touch the surface of slope.
\qfig{p2_phy_s2_para9_10_1}
\qn{9} Calculate the acceleration of the center of the ball.
\begin{qoptions}
\opt{A}{5 m/s$^2$}
\opt{B}{3.57 m/s$^2$}
\opt{C}{2.37 m/s$^2$}
\opt{D}{1.28 m/s$^2$}
\end{qoptions}
\qn{10} If the coefficient of friction is 0.2, at what value of $\tan\alpha$ will the ball slip on the rails ?
\begin{qoptions}
\opt{A}{0.14}
\opt{B}{0.28}
\opt{C}{0.63}
\opt{D}{0.38}
\end{qoptions}
\sect{Section-II(i)}
\paraq{1}{2}
A disc shaped photoelectric detector of area $\dfrac{\pi}{4}\,\mathrm{cm^2}$ generates current due to photoelectric effect. The photoelectron generation efficiency is 9\%. In the arrangement shown, $S$ is an isotropic 100 watt point monochromatic light source with $\lambda=310$nm. The focal length of the lens is 20 cm and its circular area is 4 cm$^2$. The work function is 2.0 eV.
\qfig{p2_phy_s3_para1_2_1}
\qn{1} Find the photocurrent (in mA) in the detector assuming all rays to be paraxial w.r.t. lens.
\qn{2} What is the power (in mW) of photoelectrons assuming that the average kinetic energy is half of the maximum kinetic energy?
\paraq{3}{4}
A manometer reads the pressure of a gas in a enclosure as shown in figure (a). When some of the gas is removed by a pump, the manometer reads as in (b). The liquid used in the manometers is mercury and the atmospheric pressure is 75 cm of mercury. Assume gas to be ideal.

($\rho_{\mathrm{Hg}}=14$ gm/cc, $T=27^\circ$C, $R=\dfrac{25}{3}$ J/mol k)
\qfig{p2_phy_s3_para3_4_1}
\qn{3} If the volume occupied by gas in (a) is 100 cc and the temperature of the gas remains same before and after the removal, what is the number of moles (in $\mu$ mol) of gas that is removed? Area of the manometer tube is 1cm$^2$.
\qn{4} How high (in cm) would the mercury rise in the left limb in case (b) if 14 cm of water is poured into the right limb of the manometer?
\paraq{5}{6}
Consider a simple circuit assembled from two resistors with resistance $R=1$ k$\Omega$ and two capacitors with a capacitance of $C=10$ $\mu$F, connected to the mains voltage: 200 V, 100 rad/s (fig.). Ammeter in circuit has a very small resistance. Here we have an ideal ammeter which measures the effective (rms) current through itself.
\qfig{p2_phy_s3_para5_6_1}
\qn{5} Find the ammeter reading (in A)
\qn{6} If the ammeter is replaced by an ideal hot wire Voltmeter, which measures the rms voltage, what is the voltage measured by the voltmeter (in V)?
\sect{Section-II(ii)}
\qn{7} In the centre of a long solenoid which has $n$ turns per unit length, there is a coil consisting of $N$ turns having a cross-section area $S$ (Fig.). Coil is attached to the one end of the beam balance. In the absence of current, the beam balance is in equilibrium. When a current is passed through the system, then to balance the coil, we need to add a mass $m$ on the other side as shown. The length of the right side of the lever arm is $\ell$. Determine the current (in A) in the system, if the coil and solenoid are connected in series. ($m=100\pi$gm, $N=100$, $n=1000$ turns/m, $S=100$ cm$^2$, $\ell=25$ cm.)
\qfig{p2_phy_s4_q07_1}
\qn{8} The body of mass 100 gm is suspended on a weightless spring of stiffness 500 N/m and natural length of 100 cm. Then the body is spun with a frequency of 5 Hz so that the spring describes a cone in space (Fig). Determine extension in the spring (in cm). Take : $\pi^2=10$
\qfig{p2_phy_s4_q08_1}
\qn{9} A Cart with mass $m$ moves horizontally on rails with speed $v$ (fig.). The rails go further down and smoothly transition into the new horizontal section located $H$ below. On the lower horizontal section of the rails we have another cart of mass $M$. At what initial trolley speed $v$ (in m/s) will the cart of mass $m$ be able to just rise back to the upper horizontal section again after collision with $M$? There is no friction, and the collision is elastic. ($M=9$ kg, $m=1$kg, $H=0.8$ m, $g=10$ m/s$^2$)
\qfig{p2_phy_s4_q09_1}
\raggedcolumns
\partsec{2}{CHEMISTRY}
\sect{Section-I(i)}
\qn{1} The vapour pressure of pure water at 60$^\circ$C is 150 torr. The vapour pressure of water over a solution at 60$^\circ$C containing equal numbers of moles of water and ethylene glycol ( a non volatile solute) is 50 torr. Which of the following is correct.
\begin{qoptions}
\opt{A}{The relative lowering of vapour pressure $=2/3$}
\opt{B}{The solution is ideal according to Raoult's law}
\opt{C}{The solution has positive deviation from ideal behavior}
\opt{D}{$\Delta H_{\mathrm{mix}}<0$}
\end{qoptions}
\qn{2} Potassium crystallizes in a body centred cubic lattice with edge length $a=5.0$ \AA. Then choose the correct statements in given option -
\begin{qoptions}
\opt{A}{The number of nearest neighbour to any K atom is 8}
\opt{B}{The number of next nearest neighbour to any K atom is 6}
\opt{C}{Distance between nearest neighbours is 4.33 \AA}
\opt{D}{Distance between next nearest neighbours is 5.0 \AA}
\end{qoptions}
\qn{3} A compound A was isolated as a major product formed by the reaction between a s-block metal M and excess O$_2$. A reacts with water to give only MOH, while M reacts in a controlled manner with water giving MOH and another product, B.

Select \textbf{CORRECT} for metal ``M''.
\begin{qoptions}
\opt{A}{Chloride of ``M'' gives red colour on flame.}
\opt{B}{A compound ``C'' is form on reaction with metal M \& nitrogen, produce a basic gas on hydrolysis.}
\opt{C}{Carbonate \& bicarbonates of metal ``M'' are crystalline solid.}
\opt{D}{Chloride of metal M is deliquscent \& crystalise as hydrated MCl.2H$_2$O}
\end{qoptions}
\qn{4} Metal M is a light-red metal, which is soft, malleable, and ductile. It melts at 1038$^\circ$C. Because of its positive standard electrode potential it is insoluble in hydrochloric acid and in dilute sulphuric acid, although in the presence of oxygen some dissolution might take place. Medium-concentrated nitric acid (8M) dissolves metal readily. If coloured solution MSO$_4$ obtained by dissolving metal in conc. H$_2$SO$_4$ is diluted \& treated by KI excess, precipitate of metal(I) iodide is formed, which is white, but the solution is intensely brown because of the formation of triiodide ions.

Select \textbf{CORRECT} regarding above information.
\begin{qoptions}
\opt{A}{If a clean iron nail or a blade of a penknife is immersed in a solution of a metal salt, a red deposit of metal is obtained}
\opt{B}{Metal salt solution of ``M'' shows redox reaction with Na$_2$S$_2$O$_3$}
\opt{C}{Solution obtained by reaction of metal ``M'' \& AgNO$_{3\mathrm{aq.}}$ produce colorless solution with excess NH$_4$OH}
\opt{D}{Oxide of metal M (i.e. MO) dissolves in dil. HCl but not in excess KOH}
\end{qoptions}
\qn{5} Which of the given following stability order is correct.
\begin{qoptions}
\opt{A}{\qfigi{p2_chem_s1_q05_A}}
\opt{B}{\qfigi{p2_chem_s1_q05_B}}
\opt{C}{\qfigi{p2_chem_s1_q05_C}}
\opt{D}{\qfigi{p2_chem_s1_q05_D}}
\end{qoptions}
\qn{6} In which reaction product formation takes place by Hoffmann rule ?
\begin{qoptions}
\opt{A}{\qfigi{p2_chem_s1_q06_A}}
\opt{B}{\qfigi{p2_chem_s1_q06_B}}
\opt{C}{\qfigi{p2_chem_s1_q06_C}}
\opt{D}{\qfigi{p2_chem_s1_q06_D}}
\end{qoptions}
\sect{Section-I(ii)}
\paraq{7}{8}
Radioactive tracers provide a convenient mean of measuring the small concentration encountered in determining $k_{\mathrm{sp}}$ values. Exactly 30 ml of 0.001M AgNO$_3$ solution containing radioactive silver, with an activity of 40000 counts per ml was mixed with 10 ml of 0.002 M KIO$_3$ solution and the mixture is diluted to exactly 400 ml. After equilibrium was reached, a portion of sample was filtered to remove any solid and was found to have an activity of 2000 counts per ml.
\qn{7} The $k_{\mathrm{sp}}$ of AgIO$_3$ is -
\begin{qoptions}
\opt{A}{$6\times10^{-10}$}
\opt{B}{$1.25\times10^{-9}$}
\opt{C}{$6\times10^{-12}$}
\opt{D}{$4\times10^{-10}$}
\end{qoptions}
\qn{8} Number of moles of solid AgIO$_3$ formed -
\begin{qoptions}
\opt{A}{$10^{-5}$}
\opt{B}{$10^{-4}$}
\opt{C}{$10^{-2}$}
\opt{D}{$0.4\times10^{-3}$}
\end{qoptions}
\paraq{9}{10}
The transition elements have an unparalled tendency to form coordination compounds with Lewis bases, i.e. with groups which are able to donate an electron pair. These groups are called ligands.

The reason transition elements are so good at forming complexes is that they have small, highly charged ions and have vacant low energy orbitals to accept lone pairs of electrons donated by other groups or ligands.
\qn{9} Which $\mathrm{M^{+2}}$ aquated ion/aquated complex has maximum tendency for oxidation to $\mathrm{M^{+3}_{aq.}}$.
\begin{qoptions}
\opt{A}{$\mathrm{Cr^{+2}}$}
\opt{B}{$\mathrm{Ti^{+2}}$}
\opt{C}{$\mathrm{V^{+2}}$}
\opt{D}{$\mathrm{Mn^{+2}}$}
\end{qoptions}
\qn{10} The transition metal complex(es) with zero magnetic moment, zero dipole moment and CFSE of $-2.4\ \Delta_0$ is(are)

(I) $[\mathrm{Mn(CO)_5(CH_3)}]$

(II) $[\mathrm{trans\text{-}Ni(ethylene\ diamine)_2Cl_2}]$

(III) $[\mathrm{trans\text{-}Co(CN)_4(NH_3)_2}]^{-}$

(IV) $[\mathrm{trans\text{-}Fe(CN)_4(CO)_2}]^{2-}$
\begin{qoptions}
\opt{A}{II, III}
\opt{B}{III, IV}
\opt{C}{I, III}
\opt{D}{II, III, IV}
\end{qoptions}
\sect{Section-II(i)}
\paraq{1}{2}
Resistance of a conductivity cell filled with 0.01 M HA (aq.) solution is 200 $\Omega$ at 300K. (HA is weak monoprotic acid). If cell constant is 4 cm$^{-1}$ \& $\Lambda^0_m$ (HA) is 2 $\Omega^{-1}\mathrm{m^2mol^{-1}}$. ($R=0.08$ bar/mol-K)
\qn{1} Osmotic pressure of 0.01M HA(aq.) at 300K is (in bar)-
\qn{2} Molar conductivity of 0.01M HA is $x\times10^5$ then find $x$ ($\Omega^{-1}\mathrm{cm^2mol^{-1}}$)
\paraq{3}{4}
Cuprite, Chalcosite, Galena, Pyrolusite, Heamatite, Sphelerite, Anglesite, Argentite, Chromite, Siderite, Bauxite, Cinnabar are the ores of different metals used to extract corresponding metal. Out of given ores ``X'' number of ores are sulphide ore and ``Y'' number of ores have MS type formula. An ore with formula MS with atomic mass of metal is Z, is roasted with excess O$_2$ to produce corresponding oxide.

If 1.164 gm MS produces 0.972 gm metal oxide. Which is further reduced to produce metal M.

(Atomic weight of O = 16 \& S = 32)
\qn{3} What is the value of ``X'' :-
\qn{4} What is the value of ``Z'' :-
\paraq{5}{6}
\begin{center}\small\setlength{\tabcolsep}{3pt}
\begin{tabular}{|c|l|c|p{3.4cm}|}\hline
\multicolumn{2}{|c|}{\begin{tabular}{@{}c@{}}Column-I\\(Carbohydrate)\end{tabular}} & \multicolumn{2}{c|}{\begin{tabular}{@{}c@{}}Column-II\\(Monomer unit)\end{tabular}}\\ \hline
(P) & Cellulose & (i) & $\beta$-Glucose\\ \hline
(Q) & Sucrose & (ii) & Glucopyranose \& Fructofuranose\\ \hline
(R) & Starch & (iii) & $\alpha$-Glucose\\ \hline
(S) & Lactose & (iv) & Glucose \& $\beta$-Galactose\\ \hline
\end{tabular}\end{center}
\qn{5} How many reducing sugar in given molecules

$\alpha$-Glucose, $\beta$-Glucose, Galactose, Starch
\qn{6} How many polysachride are in given molecules

$\alpha$-Glucose, $\beta$-Glucose, Galactose, Starch
\sect{Section-II(ii)}
\qn{7} The second virial coefficient of methane can be approximate by the empirical equation :

$B=a+b\cdot e^{-c/T^2}$

where $a=-0.2\ \mathrm{bar^{-1}}$, $b=0.22\ \mathrm{bar^{-1}}$, and $c=950\ \mathrm{K^2}$. What is the value of the Boyle temperature of methane (in K)? ($\ln 1.1=0.095$)
\qn{8} Write sum of the molecular weight of nitrogen containing major volatile products of given reactions.

(Given atomic weight : N = 14, O = 16, H = 1, S = 32, C = 12, Cl = 35.5 \& I = 124.00)

(i) $\mathrm{NaNO_2+KI}\xrightarrow{\ \mathrm{H^+}\ }$

(ii) $\mathrm{NaNO_2+Zn+KOH}\xrightarrow{\ \Delta\ }$

(iii) $\mathrm{NaNO_{2(aq.)}+NH^+_{4(aq.)}}\xrightarrow{\ \mathrm{Boiling}\ }$

(iv) $\mathrm{NaNO_2+}$ sulphamic acid $\xrightarrow{\ \mathrm{dil.\ HCl}\ }$
\qfullbegin
\qn{9} Compound are listed below :
\qfig{p2_chem_s4_q09_1}
\qfullsplit
Number of carboxylic acid whose decarboxylation product give +ve Iodoform test.
\raggedcolumns
\qfullend
\partsec{3}{MATHEMATICS}
\sect{Section-I(i)}
\qn{1} Let $\vec a,\ \vec b$ and $\vec c$ be three unit vectors such that $\left|\vec a+\vec b+\vec c\right|=\sqrt3$ and $\left(\vec a\times\vec b\right)\cdot\left(\vec b\times\vec c\right)+\left(\vec b\times\vec c\right)\cdot\left(\vec c\times\vec a\right)+\left(\vec c\times\vec a\right)\cdot\left(\vec a\times\vec b\right)=\lambda$ which of the following is/are \textbf{CORRECT} ?
\begin{qoptions}
\opt{A}{The maximum value of $\lambda$ is 0}
\opt{B}{If $\lambda$ is maximum, then volume of parallelopiped determined by $\vec a,\ \vec b$ and $\vec c$ is $\sqrt3$}
\opt{C}{If $\lambda$ is maximum, then value of $\bigl|(2\vec a+3\vec b+4\vec c)\cdot\allowbreak(\vec a\times\vec b+5\vec b\times\vec c+6\vec c\times\vec a)\bigr|$ is 32}
\opt{D}{None of these}
\end{qoptions}
\qn{2} If the system of equation $x+y+z=6$, $x+2y+3z=10$ and $x+2y+\lambda z=\mu$ has infinite number of solutions then which of the following is/are TRUE :
\begin{qoptions}
\opt{A}{$\dfrac{\mu-1}{\lambda}=3$}
\opt{B}{$\lambda+\mu=13$}
\opt{C}{$|\lambda-\mu|=7$}
\opt{D}{$\lambda\mu=40$}
\end{qoptions}
\qn{3} Which of the following is/are \textbf{CORRECT} ?
\begin{qoptions}
\opt{A}{$\sin1^\circ\sin2^\circ\sin3^\circ\ldots\sin90^\circ=\dfrac{3\sqrt{10}}{2^{89}}$}
\opt{B}{$\cos4^\circ\cos8^\circ\cos12^\circ\ldots\cos88^\circ=\dfrac{1}{2^{22}}$}
\opt{C}{$\cos\dfrac{\pi}{11}+\cos\dfrac{3\pi}{11}+\cos\dfrac{5\pi}{11}+\cos\dfrac{7\pi}{11}+\cos\dfrac{9\pi}{11}=\dfrac12$}
\opt{D}{None of the above option are correct}
\end{qoptions}
\qn{4} Let $P$ be any point on the curve $S=0$ such that tangents from $P$ to $x^2+y^2-2x-4y-4=0$ makes $60^\circ$ with each other and from point $Q$ perpendicular tangents are drawn to $S$ then
\begin{qoptions}
\opt{A}{Locus of $P$ is a circle of radius $2\sqrt3$}
\opt{B}{Locus of $P$ is a circle of radius 6}
\opt{C}{Locus of $Q$ is a circle of radius $3\sqrt2$}
\opt{D}{Locus of $Q$ is a circle of radius $6\sqrt2$}
\end{qoptions}
\qn{5} In $\Delta ABC$, if $BC$, $CA$ and $AB$ are three consecutive odd integers and largest angle is $120^\circ$ then
\begin{qoptions}
\opt{A}{cosine of smallest angle is $\dfrac{13}{14}$}
\opt{B}{Area of triangle $ABC$ is $\dfrac{15\sqrt3}{4}$}
\opt{C}{Inradius of $\Delta ABC$ is $\dfrac{\sqrt3}{2}$}
\opt{D}{Angle are in A.P.}
\end{qoptions}
\qn{6} Let $I_1=\displaystyle\int_1^{\sin\theta}\dfrac{x\,dx}{1+x^2}$ and $I_2=\displaystyle\int_1^{\csc\theta}\dfrac{dx}{(1+x^2)\,x}$ then which of the following is/are True :
\begin{qoptions}
\opt{A}{$\left|\dfrac{I_2}{I_1}\right|=1$}
\opt{B}{$\begin{vmatrix}I_1 & I_1^2 & I_2\\ e^{I_1+I_2} & I_2^2 & -1\\ 1 & I_1^2+I_2^2 & -1\end{vmatrix}=0$}
\opt{C}{$\dfrac{I_2}{I_1}=-2$}
\opt{D}{$(I_1+I_2)^2=1$}
\end{qoptions}
\sect{Section-I(ii)}
\paraq{7}{8}
Let $A=[a_{ij}]_{2\times2}$ be a square matrix such that $\mathrm{Tr}(A)=5$. If $f(x)=\det(A-xI)$ and $f(A)=A^2+aA+3I$ where $I$ is identify matrix of order 2 then
\qn{7} The value of $\mathrm{Tr}\left((3A^{-1})^2+(a+1)3A^{-1}+3I\right)$ is
\begin{qoptions}
\opt{A}{0}
\opt{B}{5}
\opt{C}{7}
\opt{D}{10}
\end{qoptions}
\qn{8} If $a_{ij}\in\{1,2,3,4\}$ then number of possible invertible matrices $A$ is equal to
\begin{qoptions}
\opt{A}{4}
\opt{B}{6}
\opt{C}{8}
\opt{D}{12}
\end{qoptions}
\paraq{9}{10}
Sara rolls 6 standard 6-sided cubic dice simultaneously and calculate the product of six numbers obtained.
\qn{9} Probability that product is divisible by 4
\begin{qoptions}
\opt{A}{0.9218}
\opt{B}{0.0921}
\opt{C}{0.9125}
\opt{D}{0.8215}
\end{qoptions}
\qn{10} Probability that product is divisible by 2 but not by 4
\begin{qoptions}
\opt{A}{0.921}
\opt{B}{0.0625}
\opt{C}{0.0921}
\opt{D}{0.625}
\end{qoptions}
\sect{Section-II(i)}
\paraq{1}{2}
Consider two planes $P_1:2x-3y+z=-5$ and $P_2:x-4y-14z=5$. If a line $L$ perpendicular to $P_1$ whose one point $A(-7,y_1,z_1)$ lies on plane $P_1$ and point $B(x_2,y_2,z_2)$ and $C(x_3,y_3,z_3)$ lies on plane $P_2$ as well as on line $L$.
\qn{1} Let equation of line $L$ be $\dfrac{x+7}{a}=\dfrac{y-y_1}{-b}=\dfrac{z-z_1}{c}$ where $a,b,c\in N$ then least value of $a+b+c$ is
\qn{2} Value of $\dfrac{z_3-z_2}{y_3-y_2}+\dfrac{x_3-x_2}{z_3-z_2}$ is
\paraq{3}{4}
A parabola $C$ whose focus is $S(0,0)$ and passing through $P(3,4)$. Equation of tangent at vertex of parabola is $3x+4y=25$. A chord through $S$ parallel to tangent at $P$ intersect parabola at $A$ and $B$.
\qn{3} Radius of the circumcircle of triangle formed by tangents at $A$, $B$ and chord $AB$ is
\qn{4} Area of the triangle formed by $3x+4y=25$ and tangent lines at $A$ and $B$
\paraq{5}{6}
Let $P(x)$ be a polynomial function of degree 5 such that $\displaystyle\lim_{t\to0}\left(1+\dfrac{P(t)}{t^3}\right)^{\frac1t}=e^4$ and $P(1)=9$.
\qn{5} Slope of line touching the curve $y=x^2P(x)$ at $(-1,-1)$ is
\qn{6} Sum of absolute values of $x$ where function $y=P(x)$ attains local maxima or local minima is
\sect{Section-II(ii)}
\qn{7} Let a function $f(x)$ be such that $f''(x)=f'(x)+e^x$ and $f(0)=0$, $f'(0)=1$, then value of $\ln\left(\dfrac{(f(2))^2}{4}\right)$
\qn{8} If the $(m+1)^{\text{th}}$, $25^{\text{th}}$ and $(r+1)^{\text{th}}$ term of arithmetic progression are in geometric progression and $m$, 24, $r$ are in H.P., then absolute ratio of first term of A.P. to its common difference
\qn{9} Number of distinct possible values of $[L]$, if $L=\displaystyle\lim_{x\to\pi}\dfrac{(x^2-\pi^2)^2}{1+\cos nx}\ (n\in N)$ ($[.]$ denotes greatest integer function) is -
\end{multicols}
\end{document}
